%=============================================================
% Paper 8: Krusell, Mukoyama, Rogerson & Sahin (2017), AER
%=============================================================
\chapter{Krusell, Mukoyama, Rogerson \& \c{S}ahin (2017)}

\begin{paperbox}{Gross Worker Flows over the Business Cycle}
\textbf{Authors:} Per Krusell, Toshihiko Mukoyama, Richard Rogerson,
Ayşegül Şahin\\[3pt]
\textbf{Journal:} \textit{American Economic Review} 107(11): 3447--3476\\[3pt]
\textbf{Year:} 2017 \quad\textbf{Field:} Macroeconomics; Labour economics\\[3pt]
\textbf{Classification:} Structural (heterogeneous-agent calibrated model)
\end{paperbox}

%------------------------------------------------------------
\section{Research Question}

Can a parsimonious hybrid model---combining heterogeneous-agent labour supply
with labour market frictions---simultaneously account for the cyclical
properties of all six gross worker flow rates across employment (E),
unemployment (U), and nonparticipation (N)?

%------------------------------------------------------------
\section{Motivation}

The E--U and U--E flows between employment and unemployment have been
extensively studied (Blanchard--Diamond 1990; Shimer 2012). However, flows
involving nonparticipation (N)---a state comprising $\approx$35\% of the
population---were largely ignored, even though the flows into and out of N
are as volatile as E--U flows. Two key counterintuitive facts motivate the
paper: (i) the U-to-N transition rate ($f_{UN}$) is \emph{procyclical}---
workers move from unemployment to nonparticipation \emph{more} in booms
despite the participation rate being procyclical; (ii) both pure labour
supply models (no U state) and pure search models (no participation margin)
fail to capture this pattern. The paper builds a hybrid that accounts for
all six flows simultaneously.

%------------------------------------------------------------
\section{Main Contribution}

\begin{enumerate}
  \item \textbf{Documented new facts:} All six $f_{ij}$ transition rates
    tabulated with two misclassification corrections (Abowd--Zellner;
    de$\overline{\text{NUN}}$ification). Key counterintuitive fact:
    $f_{UN}$ is \emph{procyclical}.
  \item \textbf{First model to match all six steady-state flow rates} within
    95\% confidence intervals.
  \item \textbf{Intertemporal substitution in a frictional model:} Even with
    a constant wage per efficiency unit, intertemporal labour supply
    effects emerge through the job ladder (higher $f_{EE}$ in booms
    accelerates career advancement).
  \item \textbf{Participation margin is first-order:} Flows between N and
    the other states account for $\approx$33\% of cyclical fluctuations
    in the unemployment rate.
\end{enumerate}

%------------------------------------------------------------
\section{Relationship to the Literature}

Bridges three literatures: (i) heterogeneous-agent labour supply
(Lucas--Rapping 1969; Chang--Kim 2006); (ii) search-matching
(Mortensen--Pissarides 1994; Shimer 2012); (iii) gross worker flows with
participation (Elsby, Hobijn, \c{S}ahin 2015; Shimer 2013). Distinguished
from existing hybrid models by focusing on \emph{gross worker flows} (not
stocks alone) and by introducing UI, on-the-job search, and heterogeneous
match quality for realistic calibration.

%------------------------------------------------------------
\section{Theoretical Framework and Economic Mechanism}

\subsection*{Individual Preferences}

\begin{equation}
  \mathbb{E}_t \sum_{t=0}^{\infty}\beta^t\bigl[\log(c_t) - \alpha e_t
    - \gamma s_t\bigr],
  \label{eq:kmrs_prefs}
\end{equation}
where $c_t\geq 0$ is consumption, $e_t\in\{0,1\}$ employment status,
$s_t\in\{0,1\}$ active search indicator, $\alpha>0$ (disutility of work),
$\gamma>0$ (disutility of search).

\subsection*{Idiosyncratic Productivity and Match Quality}

$\log z_{t+1} = \rho_z\log z_t + \varepsilon_{t+1}$ with
$\varepsilon_{t+1}\sim N(0,\sigma_\varepsilon^2)$.

Each employment opportunity comes with an i.i.d.\ match quality draw
$q\sim\text{LogNormal}(0,\sigma_q^2)$; employed workers can accept better
matches via on-the-job search ($\lambda_e$).

\subsection*{Labour Market Frictions}

Four friction parameters characterise the aggregate environment
$\Lambda=(w,r,\lambda_u,\lambda_n,\lambda_e,\sigma)$:
\begin{itemize}
  \item $\lambda_u$: offer arrival rate for active searchers (unemployed)
  \item $\lambda_n$: offer arrival rate for passive searchers (out of LF)
  \item $\lambda_e$: outside offer rate for employed workers
  \item $\sigma$: job separation rate
\end{itemize}

\subsection*{Bellman Equations (key)}

For a worker with an employment opportunity (value $W$ if accepted, $J$ if
rejected):
\begin{multline}
  W(a,z,q,I^B,\Lambda) = \max_{c\geq 0,a'\geq 0}\Bigl\{\ln c - \alpha \\
    + \beta\,\mathbb{E}\bigl[(1-\sigma-\lambda_e)V(a',z',q,\gamma',0,\Lambda')
    + \lambda_e V(a',z',\max\{q,q'\},\gamma',0,\Lambda') \\
    + \sigma(1-\lambda_u)J(a',z',\gamma',1,\Lambda')
    + \sigma\lambda_u V(a',z',q',\gamma',1,\Lambda')\bigr]\Bigr\}
\end{multline}
subject to $c+a'=(1+r)a+(1-\tau)wzq+T$.

For an unemployed worker who chooses active search:
\begin{equation}
  U(a,z,\gamma,I^B,\Lambda) = \max_{c,a'}\bigl\{\ln c - \gamma + \beta\,\mathbb{E}\bigl[\lambda_u V + (1-\lambda_u)J\bigr]\bigr\}
\end{equation}
subject to $c+a'=(1+r)a+(1-\tau)I^B b(z)+T$.

\subsection*{Mechanism for Procyclical $f_{UN}$}

In booms, strongly attached workers find jobs quickly, leaving the
unemployment pool populated by workers near the U--N indifference margin.
These workers are more likely to transition to N, generating the
\emph{counterintuitive} procyclical $f_{UN}$.

%------------------------------------------------------------
\section{Data}

\begin{itemize}
  \item \textbf{CPS gross worker flows:} Monthly CPS matched surveys,
    1978:Q1--2012:Q3; six $f_{ij}$ rates computed (Table~2).
  \item \textbf{Misclassification corrections:} (1) Abowd--Zellner (1985)
    reinterview-based; (2) de$\overline{\text{NUN}}$ification (recodes
    U$\to$N$\to$U sequences as N). Both substantially reduce flows involving N.
  \item \textbf{SIPP:} Validates model predictions for flow rates conditional
    on wealth.
  \item \textbf{Calibration targets:} LFPR (0.66), unemployment rate (0.068),
    $f_{EU}$ (0.014), $f_{NE}$ (0.022); J2J rate (2.2\%/month); avg wage
    gain on J2J (3.3\%).
\end{itemize}

%------------------------------------------------------------
\section{Empirical Strategy}

\begin{enumerate}
  \item Tabulate cyclical properties of all six gross flow rates using
    HP-filtered ($\lambda=1600$) quarterly series (Table~3).
  \item Calibrate the model's steady state to match average flow rates
    (Table~4--5).
  \item Subject calibrated model to aggregate shocks (drawn from cyclical
    variation in job-finding and separation rates); compare model-generated
    cyclical moments to data (Table~6).
\end{enumerate}

%------------------------------------------------------------
\section{Identification Strategy}

Calibration, not econometric identification. Parameters are set through:
(a)~external values ($\rho_z=0.996$, $\sigma_\varepsilon=0.096$ from wage
panel studies; $\tau=0.30$; UI parameters from programme data);
(b)~steady-state moment matching (Table~4); (c)~model fixed-point conditions.
Aggregate shocks are measured from data on cyclical variation in job-finding
and separation rates.

%------------------------------------------------------------
\section{Main Results}

\subsection*{Key Cyclical Facts (Table~3, AZ-corrected)}

\begin{center}
\begin{tabular}{lccc}
\toprule
\textbf{Flow} & \textbf{Std dev} & $\boldsymbol{\rho(x,Y)}$ &
  \textbf{Pattern} \\
\midrule
$f_{EU}$ & 0.089 & $-0.63$ & Countercyclical \\
$f_{EN}$ & 0.083 & $+0.43$ & \textbf{Procyclical} (counterintuitive) \\
$f_{UE}$ & 0.088 & $+0.76$ & Procyclical \\
$f_{UN}$ & 0.106 & $+0.61$ & \textbf{Procyclical} (counterintuitive) \\
$f_{NE}$ & 0.103 & $+0.52$ & Procyclical \\
$f_{NU}$ & 0.072 & $-0.23$ & Weakly countercyclical \\
\bottomrule
\end{tabular}
\end{center}

Note: $f_{UN}$ procyclical means workers leave unemployment for N
\emph{more} in booms---counterintuitive from a discouraged-worker perspective.

\subsection*{Calibration Achievement (Table~5)}

The model matches all six average flow rates within 95\% CIs---the first
structural model to achieve this. Previous models could not match N-adjacent
flows.

\subsection*{Cyclical Performance (Table~6)}

Model does a good job matching most cyclical moments; flows between
participation and nonparticipation account for $\approx$33\% of unemployment
rate fluctuations.

%------------------------------------------------------------
\section{Economic Magnitude and Interpretation}

The participation margin is far more important for cyclical unemployment than
previously recognised: one-third of the rise in the unemployment rate during
recessions reflects unemployed workers failing to exit to employment
($\downarrow f_{UE}$, $\uparrow f_{EU}$), but also the participation margin
($f_{UN}$, $f_{NE}$). The procyclical $f_{UN}$ implies the standard
``discouraged worker'' narrative (workers give up searching in recessions) is
a level effect, not a rate effect: the stock of discouraged workers rises in
recessions not because the rate $f_{UN}$ increases but because the stock of
unemployed is larger.

%------------------------------------------------------------
\section{Mechanisms}

Two forces nearly cancel to produce the \emph{weak} procyclicality of the
participation rate:
\begin{itemize}
  \item \textbf{Wealth effect:} Higher assets in booms reduce desire to work
    (reduces participation).
  \item \textbf{Intertemporal substitution:} Higher $f_{EE}$ (J2J) in booms
    accelerates movement up the job-quality ladder $\Rightarrow$ higher
    return to participation now (increases participation).
\end{itemize}
Near-cancellation $\Rightarrow$ LFPR fluctuates little ($\sigma=0.0026$)
compared to the unemployment rate ($\sigma=0.117$).

Procyclical $f_{UN}$: in booms, the unemployment pool is disproportionately
composed of workers close to the U--N indifference margin; these workers
transition to N at a higher rate, generating the counterintuitive procyclicality.

%------------------------------------------------------------
\section{Robustness Checks}

\begin{itemize}
  \item Two misclassification corrections (Abowd--Zellner and
    de$\overline{\text{NUN}}$ification) give quantitatively similar cyclical
    properties; both shown in Table~3.
  \item Model evaluated against both correction methods in Table~6.
  \item SIPP wealth-conditional flow data used to validate model heterogeneity
    predictions.
  \item An earlier GE version (Krusell et al.\ 2012) gives similar results,
    confirming the PE assumption is not driving findings.
\end{itemize}

%------------------------------------------------------------
\section{Limitations}

\begin{enumerate}
  \item \textbf{Partial equilibrium:} Prices and aggregate frictions are
    exogenous; wages and vacancy posting not endogenised.
  \item \textbf{Friction shocks only:} The source of cyclical friction
    variation is not connected to TFP, monetary policy, or aggregate demand.
  \item \textbf{No age, gender, or education heterogeneity:} All empirically
    important for N dynamics, especially for participation decisions.
  \item \textbf{Quantitative mismatch:} Table~6 shows the model does not
    fully replicate all six cyclical moments simultaneously.
\end{enumerate}

%------------------------------------------------------------
\section{Policy Implications}

\begin{itemize}
  \item Policies targeting unemployment (UI, ALMP) affect flows from N
    as well---the participation margin is a first-order policy variable.
  \item The intertemporal substitution channel implies aggregate demand
    stimulus that increases J2J flows also raises participation---a positive
    feedback loop.
  \item Procyclical $f_{UN}$: UI recipients are more likely to stop searching
    and exit to N in \emph{booms} (after finding a job, they stop claiming),
    consistent with UI moral hazard being cyclically variable.
\end{itemize}

%------------------------------------------------------------
\section{Overall Assessment}

A technically accomplished paper that fills an important gap in macro-labour
economics. The empirical documentation of all six gross flow rates with careful
misclassification corrections is valuable in itself. The calibration
achievement---matching all six average flow rates within CIs---is a genuine
advance. The conceptual contributions (participation fluctuations are first-
order; intertemporal substitution operates through the job ladder) are
important and influential.

\textbf{Quality: Very good. Strong data work, careful calibration,
novel model insights.}

%------------------------------------------------------------
\section{Relevance to My Research}

Directly relevant for RDC work on Canadian labour market flows. The E/U/N
three-state framework is applicable to matched LFS data or LEED/LWF. Testing
whether Canadian gross flows show the same counterintuitive cyclical patterns
($f_{UN}$ procyclical; large N-adjacent volatility) would be a valuable
contribution, especially distinguishing resource-intensive from manufacturing
provinces.

%------------------------------------------------------------
\section{Potential Research Extensions}

\begin{itemize}
  \item Implement the gross-flow framework for Canada using matched monthly
    LFS microdata; test whether the counterintuitive procyclical $f_{UN}$
    holds in Canada and differs across provinces.
  \item Extend the model with sectoral heterogeneity (resource vs.\
    manufacturing vs.\ services) to understand how commodity cycles drive
    differential E/U/N flows in Canada.
\end{itemize}

%------------------------------------------------------------
\section{Research Gaps}

\begin{itemize}
  \item The cyclical source of friction variation is unexplained; integrating
    a production/vacancy-posting side (Mortensen--Pissarides structure) would
    close this gap.
  \item No treatment of long-run participation trends (declining since 2000);
    a low-frequency extension would be valuable.
\end{itemize}

%------------------------------------------------------------
\section{Methodological Lessons}

\begin{enumerate}
  \item The Abowd--Zellner misclassification correction is essential for
    flows involving N; spurious U--N and N--U flows dramatically inflate
    measured N-adjacent volatility.
  \item Matching all six \emph{average} flow rates (not just aggregate stocks)
    is a far more stringent model test; passing the flow test automatically
    passes the stock test, but not vice versa.
  \item Labour market stocks are a less stringent check on models than gross
    flows---when gross flows are available, they should be used as the primary
    calibration target.
\end{enumerate}

%------------------------------------------------------------
\section*{Five-Sentence Summary}

Krusell, Mukoyama, Rogerson, and \c{S}ahin (2017) document that gross worker
flows into and out of nonparticipation are as volatile over the business cycle
as the much-studied employment--unemployment flows, and that the U-to-N
transition rate is procyclical---counterintuitive but explained by the changing
composition of the unemployment pool in booms. They build a parsimonious hybrid
model combining heterogeneous-agent labour supply with search frictions, the
first structural model to match all six average gross flow rates within their
95\% confidence intervals. Despite a constant wage per efficiency unit,
intertemporal substitution operates through the job-quality ladder---higher
job-to-job flows in booms increase career returns, generating a procyclical
incentive to participate. Flows between participation and nonparticipation
account for approximately one-third of cyclical fluctuations in the
unemployment rate---far larger than previously recognised. The paper
establishes both empirically and theoretically that the participation margin
is a first-order dimension of aggregate labour market dynamics that cannot be
omitted from business cycle models of unemployment.

\vspace{6pt}
\begin{quickref}
\begin{tabularx}{\linewidth}{@{}lX@{}}
\textbf{Key contribution} &
  First model to match all six steady-state E/U/N gross flow rates within
  CIs; establishes the counterintuitive procyclicality of $f_{UN}$; documents
  that the N margin accounts for $\approx$33\% of unemployment fluctuations.
  \\[4pt]
\textbf{Key weakness} &
  Partial equilibrium; aggregate friction shocks taken as given rather than
  derived; model does not fully match all cyclical moments quantitatively;
  no demographic heterogeneity. \\[4pt]
\textbf{Most interesting gap} &
  The source of cyclical variation in friction parameters is left unexplained;
  a production side with vacancy posting and TFP shocks would complete the
  model. \\[4pt]
\textbf{Publishable extension} &
  Implement the gross-flow framework for Canada using matched monthly LFS data;
  test whether the counterintuitive procyclical $f_{UN}$ holds in Canada and
  differs between resource-intensive and manufacturing provinces.
\end{tabularx}
\end{quickref}
